%! AP Statistics — Unit 3, Lesson 3.6 — Homework KEY
\documentclass[10pt]{article}
\usepackage{apstats-article}
\usepackage{apstats-key}
\usepackage{grading}
\gradingassignment{apstats-u03-l06-hw}
\gradingcourse{advanced_statistics}

\begin{document}

\pageheader{Unit 3, Lesson 3.6}{Homework — KEY}
\namedateperiod

\vspace{0.1in}

\begin{practicebox}
\small
\textbf{1.}
\begin{enumerate}[label=(\alph*), leftmargin=2em, itemsep=5pt]
  \item \ans{Block design (blocking by gender)}\gradedkey{\gradedblank[id=hw1a, type=short, points=1, answer={randomized block design}, check=text, display={\bfseries block design (by gender)}]{2.2in}{0.2in}}
  \item \ans{Gender is a known variable that may affect anxiety response to mindfulness training. Blocking removes this variability, giving a more precise estimate of the treatment effect. With 20 men and 20 women, random assignment within blocks is straightforward.}\gradedkeyline{\gradedblank[id=hw1b, type=rubric, points=2, display={Gender may affect the response to mindfulness; blocking removes that variability, so the estimate is more precise.}]{5.2in}{\gradedlinesheight{2}}}
  \item \begin{teachernote}Diagram: 40 students $\to$ Block 1 (men, 20) / Block 2 (women, 20); within each block, randomly assign 10 to mindfulness, 10 to control.\end{teachernote}\gradedkeyline{\gradedblank[id=hw1c, type=rubric, points=3, display={40 students $\to$ Block 1 (20 men), Block 2 (20 women); within each block randomly assign 10 to mindfulness, 10 to control; compare anxiety.}]{6.1in}{1.2in}}
\end{enumerate}

\vspace{6pt}
\textbf{2.}
\begin{enumerate}[label=(\alph*), leftmargin=2em, itemsep=5pt]
  \item \ans{Matched pairs design}\gradedkey{\gradedblank[id=hw2a, type=short, points=1, answer={matched pairs design}, check=text]{2.2in}{0.2in}}
  \item \ans{Each patient controls for their own individual factors (diet, saliva chemistry, brushing habits). By applying both brands to the same patient (one side each), all individual differences are controlled, giving maximum precision.}\gradedkeyline{\gradedblank[id=hw2b, type=rubric, points=2, display={Each patient gets both brands, so diet, saliva, and brushing habits are controlled.}]{3.3in}{\gradedlinesheight{2}}}
\end{enumerate}

\vspace{6pt}
\textbf{3.}
\begin{enumerate}[label=(\alph*), leftmargin=2em, itemsep=5pt]
  \item \ans{CRD (completely randomized design)}\gradedkey{\gradedblank[id=hw3a, type=short, points=1, answer={completely randomized design}, check=text, display={\bfseries CRD}]{2.2in}{0.2in}}
  \item \ans{All seedlings are from the same variety grown under the same conditions — they are homogeneous. There is no known nuisance variable to block on, so a CRD is the simplest and most appropriate design.}\gradedkeyline{\gradedblank[id=hw3b, type=rubric, points=2, display={Seedlings are homogeneous (same variety, same conditions): no nuisance variable to block on.}]{3.8in}{\gradedlinesheight{2}}}
\end{enumerate}

\vspace{6pt}
\textbf{4.} \ans{A CRD is preferable when the blocking variable is not measurable before the study begins, when the variable's effect on the response is very small (blocking would gain little precision), or when the added logistical complexity of blocking outweighs the benefit. For example, if the nuisance variable is only weakly related to the response, blocking may not meaningfully reduce variability.}\gradedkeyline{\gradedblank[id=hw4, type=rubric, points=2, display={When the variable can't be measured beforehand, is only weakly related to the response, or blocking's complexity outweighs the gain in precision.}]{6.1in}{\gradedlinesheight{3}}}
\end{practicebox}

\end{document}
